\chapter{留数定理及其在实积分中的应用}

\problem{8.1}{求下列函数在孤立奇点的留数（其中 $n,m$ 为自然数），并求出函数孤立奇点的留数之和。
\begin{enumerate}
  \item $f(z)=\dfrac{1}{z(z-1)^2}$ \qquad
  \item $f(z)=z^n\sin\dfrac{1}{z}$
  \item $f(z)=\dfrac{\mathrm{e}^z-1}{z^5}$ \qquad
  \item $f(z)=\dfrac{\mathrm{e}^{z^2}}{(z-1)^2}$
  \item $f(z)=\dfrac{\mathrm{e}^z}{(z-1)^2}$ \qquad
  \item $f(z)=\mathrm{e}^{\frac{1}{1-z}}$
  \item $f(z)=\dfrac{z^2}{\cos z-1}$ \qquad
  \item $f(z)=z^3\cos\dfrac{1}{z-2}$
  \item $f(z)=\dfrac{1}{(z-\alpha)(z-\beta)^m}$ \qquad
  \item $f(z)=\dfrac{1}{1+z^{2m}}$
  \item $f(z)=\dfrac{z^{2m}}{1+z^m}$ \qquad
  \item $f(z)=\dfrac{\tg z}{1-\mathrm{e}^z}$
\end{enumerate}}
\answerarea

\problem{8.2}{计算下列回路积分
\begin{enumerate}
  \item $\displaystyle\oint_C \frac{1}{(z^2+1)(z-1)^2}\,\mathrm{d}z \qquad C:\ |z-1|=1$
  \item $\displaystyle\oint_C \frac{1}{(z^2+1)(z-1)^2}\,\mathrm{d}z \qquad C:\ |z-\mathrm{i}|=1$
  \item $\displaystyle\oint_C \frac{1}{(z^2+1)(z-1)^2}\,\mathrm{d}z \qquad C:\ |z|=2$
  \item $\displaystyle\oint_{|z|=2} \frac{1}{(z-3)(z^4-1)}\,\mathrm{d}z$ \qquad
  \item $\displaystyle\oint_{|z|=3} \cot^3 z\,\mathrm{d}z$
  \item $\displaystyle\oint_{|z|=1} \frac{z\sin z}{(1-\mathrm{e}^z)^3}\,\mathrm{d}z$ \qquad
  \item $\displaystyle\oint_{|z|=1} z^n\mathrm{e}^{\frac{1}{z}}\,\mathrm{d}z$
\end{enumerate}}
\answerarea
\clearpage

\problem{8.3}{计算下列积分
\begin{enumerate}
  \item $\displaystyle\int_0^{2\pi} \frac{1}{2+\cos\theta}\,\mathrm{d}\theta$ \qquad
  \item $\displaystyle\int_0^{2\pi} \cos^{2n}\theta\,\mathrm{d}\theta$
  \item $\displaystyle\int_0^{2\pi} \frac{1}{1+\sin^2\theta}\,\mathrm{d}\theta$ \qquad
  \item $\displaystyle\int_0^{2\pi} \frac{1}{1-2p\cos\theta+p^2}\,\mathrm{d}\theta,\quad p>0$
\end{enumerate}}
\answerarea

\problem{8.4}{计算下列积分
\begin{enumerate}
  \item $\displaystyle\int_0^\infty \frac{1}{x^4+1}\,\mathrm{d}x$ \qquad
  \item $\displaystyle\int_0^\infty \frac{x^2}{(x^2+a^2)^2}\,\mathrm{d}x$
  \item $\displaystyle\int_0^\infty \mathrm{e}^{-ax^2}\cos bx\,\mathrm{d}x,\quad a>0,\,b\in R$
  \item $\displaystyle\int_0^\infty \frac{(\cos 2x)^2}{(x^2+1)^2(x^4+4)}\,\mathrm{d}x$ \qquad
  \item $\displaystyle\int_{-\infty}^\infty \frac{x^2}{(x^2+1)^2(x^2+2x+2)}\,\mathrm{d}x$
  \item $\displaystyle\int_{-\infty}^\infty \frac{\mathrm{e}^{ax}}{1+\mathrm{e}^x}\,\mathrm{d}x,\quad 0<a<1$
\end{enumerate}}
\answerarea
\clearpage

\problem{8.5}{计算下列积分
\begin{enumerate}
  \item $\displaystyle\int_0^\infty \frac{\cos mx}{1+x^4}\,\mathrm{d}x,\quad m>0$ \qquad
  \item $\displaystyle\int_{-\infty}^\infty \frac{x\sin x}{x^2+a^2}\,\mathrm{d}x,\quad a>0$
  \item $\displaystyle\int_0^\infty \frac{\sin^2 x}{x^2}\,\mathrm{d}x$ \qquad
  \item $\displaystyle\int_{-\infty}^\infty \frac{x\cos x}{x^2-2x+10}\,\mathrm{d}x$
  \item $\displaystyle\int_0^\infty \frac{\sin mx}{x(x^2+a^2)}\,\mathrm{d}x,\quad a>0,\,m\text{ 为正整数}$ \qquad
  \item $\displaystyle\int_0^\infty \frac{1}{1+x^3}\,\mathrm{d}x$
\end{enumerate}}
\answerarea

\problem{8.6}{计算下列积分
\begin{enumerate}
  \item $\displaystyle P\int_{-\infty}^\infty \frac{1}{x^3-1}\,\mathrm{d}x$ \qquad
  \item $\displaystyle P\int_{-\infty}^\infty \frac{\mathrm{e}^{\mathrm{i}kx}-1}{x^2}\,\mathrm{d}x,\quad k>0$
  \item $\displaystyle P\int_{-\infty}^\infty \frac{1}{x(x-1)(x-2)}\,\mathrm{d}x$ \qquad
  \item $\displaystyle P\int_0^\infty \frac{\sin(x+a)\sin(x-a)}{x^2-a^2}\,\mathrm{d}x$
\end{enumerate}}
\answerarea
\clearpage

\problem{8.7}{在下列情况下计算积分 $\displaystyle\int_{-\infty}^\infty \frac{1}{x^2-a^2}\,\mathrm{d}x,\quad a>0$
\begin{enumerate}
  \item 让 $a\to a+\mathrm{i}\varepsilon,\quad \varepsilon\to 0,\,\varepsilon>0$。
  \item 让 $a\to a-\mathrm{i}\varepsilon,\quad \varepsilon\to 0,\,\varepsilon>0$。
  \item 直接计算该积分的 Cauchy 主值积分。
\end{enumerate}}
\answerarea
\clearpage
